\documentclass[11pt]{article}
\usepackage[spanish]{babel}
\usepackage{amsmath}
\usepackage{geometry}
\usepackage{graphicx}
\usepackage{setspace}
\geometry{margin=1.5cm}

\title{Módulo 1: Fundamentos de Stress Testing en Sistemas de Pensiones}
\author{MACROPOL}
\date{}

\begin{document}

\maketitle


\section{Introducción}

El stress testing se ha consolidado como una herramienta central en la gestión de riesgos y la regulación financiera, especialmente después de la crisis financiera global. Su objetivo es evaluar la resiliencia de sistemas financieros ante escenarios adversos plausibles.

En el contexto de los sistemas de pensiones de capitalización individual, el stress testing permite analizar cómo choques macroeconómicos, financieros y demográficos afectan los niveles de pensión futura y la sostenibilidad del sistema.

Esta nota técnica tiene como propósito presentar un marco conceptual y metodológico para la aplicación de ejercicios de stress testing en sistemas previsionales, integrando dimensiones macroeconómicas, financieras y actuariales. En particular, la nota describe cómo construir escenarios macro-financieros, cómo modelar su transmisión hacia los portafolios de inversión y las variables demográficas, y cómo evaluar su impacto sobre indicadores clave como la tasa de reemplazo y la suficiencia de las pensiones.

En este marco, el objetivo de este curso es capacitar a los participantes en el diseño e implementación de ejercicios de stress testing aplicados a sistemas previsionales de capitalización individual, integrando de manera coherente los riesgos macroeconómicos, financieros y actuariales. El curso busca desarrollar capacidades para construir escenarios adversos, modelar su transmisión hacia los portafolios de inversión y variables demográficas, y evaluar su impacto sobre la sostenibilidad del sistema y el bienestar de los afiliados.

\section{¿En qué consiste el Stress Testing?}

Un stress test es una simulación que evalúa el comportamiento de un sistema bajo escenarios extremos pero plausibles. Formalmente, se puede representar como:

\begin{equation}
Resultados = f(Escenario, Modelo, Datos)
\end{equation}

Donde:

\begin{itemize}
\item Escenario: trayectoria de variables macroeconómicas bajo condiciones adversas
\item Modelo: conjunto de relaciones económicas y financieras
\item Datos: información del sistema (afiliados, portafolio, demografía)
\end{itemize}

Se distinguen dos tipos de escenarios:

\begin{itemize}
\item Baseline: evolución esperada o escenario central
\item Adverso: condiciones de estrés o escenario estresado
\end{itemize}

\section{Flujo de trabajo del Stress Testing}

El desarrollo de un ejercicio de stress testing sigue una secuencia estructurada:

\begin{enumerate}
\item Definición del objetivo
\item Construcción de escenarios macroeconómicos
\item Identificación de factores de riesgo
\item Modelación del portafolio
\item Modelación de afiliados
\item Cálculo de pensiones
\item Evaluación de impacto
\end{enumerate}

Este flujo transforma un shock macroeconómico en un resultado relevante para política pública.

\section{Información requerida}

La implementación de un stress test requiere cuatro bloques de información, resumidos en la siguiente tabla:

\begin{table}[h]
\centering
\begin{tabular}{|l|l|}
\hline
\textbf{Bloque de información} & \textbf{Variables clave} \\
\hline

Datos Macroeconómicos &
\begin{tabular}[c]{@{}l@{}}
- Retornos de activos \\
- Inflación \\
- Crecimiento salarial
\end{tabular} \\
\hline

Datos Financieros &
\begin{tabular}[c]{@{}l@{}}
- Composición del portafolio \\
- Volatilidad y correlaciones
\end{tabular} \\
\hline

Datos de Afiliados &
\begin{tabular}[c]{@{}l@{}}
- Edad \\
- Salario \\
- Saldo acumulado \\
- Densidad de cotización
\end{tabular} \\
\hline

Datos Actuariales &
\begin{tabular}[c]{@{}l@{}}
- Tablas de mortalidad \\
- Edad de retiro
\end{tabular} \\
\hline

\end{tabular}
\caption{Información requerida para la implementación de un stress test previsional}
\end{table}

\section{Herramientas analíticas}

El desarrollo de un stress test requiere dominar un conjunto de herramientas analíticas, resumidas en la siguiente tabla:

\begin{table}[h]
\centering
\begin{tabular}{|l|l|}
\hline
\textbf{Tipo de Herramienta} & \textbf{Modelos / Técnicas} \\
\hline

Modelos Macroeconómicos &
\begin{tabular}[c]{@{}l@{}}
- ARIMA \\
- VAR
\end{tabular} \\
\hline

Modelos Financieros &
\begin{tabular}[c]{@{}l@{}}
- CAPM \\
- GARCH
\end{tabular} \\
\hline

Simulación &
\begin{tabular}[c]{@{}l@{}}
- Monte Carlo
\end{tabular} \\
\hline

Modelos Actuariales &
\begin{tabular}[c]{@{}l@{}}
- Cálculo de anualidades \\
- Modelos de supervivencia
\end{tabular} \\
\hline

Integración &
\begin{tabular}[c]{@{}l@{}}
- Simulación dinámica integrada
\end{tabular} \\
\hline

\end{tabular}
\caption{Herramientas analíticas para el desarrollo de stress testing previsional}
\end{table}
\section{Rol en Sistemas de Pensiones}

El stress testing permite:

\begin{itemize}
\item Evaluar la suficiencia de las pensiones
\item Analizar la sostenibilidad del sistema
\item Informar decisiones regulatorias
\end{itemize}

A diferencia del sector bancario, el foco no es la solvencia inmediata, sino los resultados de largo plazo.

\section{Particularidades en Pensiones}

Los sistemas previsionales presentan características específicas:

\begin{itemize}
\item Horizonte de largo plazo
\item Importancia de la longevidad
\item Dependencia de trayectorias laborales
\end{itemize}

El resultado relevante es la pensión futura, no el capital acumulado.

\section{Conclusión}

El stress testing en pensiones es un proceso integrado que combina modelos macroeconómicos, financieros y actuariales para evaluar riesgos sistémicos y apoyar la formulación de políticas públicas.

\end{document}